\documentclass[10pt,a4paper]{article}

% ---------- Packages ----------
\usepackage[utf8]{inputenc}
\usepackage[T1]{fontenc}
\usepackage[english]{babel}

\usepackage[
    a4paper,
    margin=1.8cm,
    top=1.2cm,
    bottom=1.4cm,
    includehead,
    headheight=1.8cm,
    headsep=0.4cm
]{geometry}

\usepackage{graphicx}
\usepackage[table]{xcolor}
\usepackage{titlesec}
\usepackage{enumitem}
\usepackage{multicol}
\usepackage{fancyhdr}
\usepackage{array}
\usepackage{tabularx}
\usepackage{xurl}
\usepackage{hyperref}

% ---------- Colors ----------
\definecolor{mainblue}{RGB}{25,55,110}
\definecolor{lightgray}{RGB}{245,245,245}

% ---------- Links ----------
\hypersetup{colorlinks=true, linkcolor=mainblue, urlcolor=mainblue, citecolor=mainblue}

% ---------- Header ----------
\pagestyle{fancy}
\fancyhf{}
\renewcommand{\headrulewidth}{0.4pt}
\renewcommand{\footrulewidth}{0pt}
\fancyhead[L]{\begin{minipage}[c]{0.30\textwidth}\raggedright\includegraphics[width=3.2cm]{\detokenize{uzi_logo.png}}\end{minipage}}
\fancyhead[C]{\begin{minipage}[c]{0.36\textwidth}\centering\textbf{\small Thesis Proposal}\\[-0.1em]\scriptsize Department of Informatics\end{minipage}}
\fancyhead[R]{\begin{minipage}[c]{0.30\textwidth}\raggedleft\includegraphics[width=3.0cm]{\detokenize{everlogo.png}}\end{minipage}}
\fancyfoot[C]{\thepage}

% ---------- Section Formatting ----------
\titleformat{\section}{\color{mainblue}\normalfont\normalsize\bfseries}{}{0em}{}
\titlespacing*{\section}{0pt}{0.55em}{0.25em}

% ---------- Compact Lists ----------
\setlist[itemize]{noitemsep, topsep=1pt, leftmargin=1.2em}
\setlist[enumerate]{noitemsep, topsep=1pt, leftmargin=1.4em}

% ---------- Tables ----------
\renewcommand{\arraystretch}{1.05}
\setlength{\tabcolsep}{5pt}
\setlength{\parskip}{2pt}

\begin{document}

% ---------- Title ----------
\begin{center}
    {\large \textbf{Retrieval-Augmented Generation for Multilingual Industrial\\[0.1em]Engineering Documentation: Evaluation and Abstention}}\\[0.2em]
    {\small Industry Collaboration with Everllence SE (formerly MAN Energy Solutions)}
\end{center}

\vspace{0.2em}

\noindent
{\small
\begin{tabularx}{\textwidth}{>{\bfseries}l X >{\bfseries}l X}
Student: & Necati Furkan Colak & Student ID: & 24746323 \\
Program: & MSc Computer Science & Semester: & Spring 2026 \\
Supervisor: & [Supervisor Name] & Co-Supervisor: & [Everllence Contact] \\
Email: & necatifurkan.colak@uzh.ch & Date: & \today \\
\end{tabularx}}

\vspace{0.3em}
\hrule
\vspace{0.3em}

\section{Abstract}

Engineering organisations hold large multilingual corpora of manuals, service reports, and norms; finding the correct procedure or parameter in them is a daily time cost. Retrieval-augmented generation (RAG) \cite{lewis2020rag} enables natural-language access to such corpora. This thesis builds and evaluates a RAG system with Everllence as industrial partner and adds one focused scientific increment: \emph{abstention}. In safety-critical maintenance, a fluent but wrong torque value is worse than no answer, so the system must decline when evidence is insufficient. The work builds a modular pipeline (table-aware ingestion, hybrid retrieval \cite{robertson2009bm25, karpukhin2020dpr}, Azure OpenAI + one open-weight generator), constructs an expert-validated evaluation set \emph{including unanswerable questions}, and measures correctness, groundedness (RAGAS/ARES \cite{es2023ragas, saadfalcon2024ares}), and the abstention trade-off, validating automatic metrics against expert judgement.

\section{Motivation and Gap}

Everllence engineers consult documentation spanning product generations, languages (German/English), and formats \cite{everllence2025}; search is keyword-based and retrieval expertise is concentrated in senior staff. Building a RAG demo is easy --- knowing whether it can be trusted is not: technical facts live in tables and variant-specific exceptions, hallucination is a documented failure mode \cite{ji2023survey}, and its cost here is asymmetric. Automatic evaluation frameworks \cite{es2023ragas, saadfalcon2024ares} and benchmarks measuring ``negative rejection'' \cite{chen2024benchmark} exist, but have been validated mainly on open-domain corpora. Whether their scores track \emph{expert} judgement on dense engineering documentation --- and whether prompting-based abstention gives a usable wrong-answer/coverage trade-off --- are open, answerable questions.

\section{Feasibility}

\begin{itemize}
    \item Bounded workload: one corpus pair, $\sim$4 retrieval configurations, 2 generators, prompt-level abstention only --- the scientific effort concentrates on evaluation design, not risky engineering.
    \item Mature tooling (open-source ingestion/retrieval frameworks, RAGAS); Everllence's Azure OpenAI licences remove infrastructure risk.
    \item Data risk hedged: methodology is developed on a public technical-documentation twin corpus (open engine manuals, classification-society rules); company documents enrich the study but are not blocking.
    \item Expert effort bounded: two batched workshops (question authoring; blinded rating).
\end{itemize}

\section{Research Questions}

\begin{enumerate}
    \item \textbf{RQ1:} How accurately and groundedly does a well-configured RAG pipeline answer realistic engineering questions over multilingual technical documentation, and which retrieval configuration (chunking; BM25/dense/hybrid) performs best?
    \item \textbf{RQ2:} Can prompting-based abstention reliably reject unanswerable questions, and at what false-refusal cost on answerable ones?
    \item \textbf{RQ3:} How well do automatic RAG metrics agree with domain-expert judgement on this corpus --- can routine evaluation be automated for regression testing?
\end{enumerate}

\section{Methodology}

\begin{itemize}
    \item \textbf{Corpora:} public twin corpus (German/English manuals, rules, standards) for reproducibility + curated non-confidential Everllence set in the company's Azure tenant; table-preserving, structure-aware ingestion.
    \item \textbf{Evaluation set:} 150--200 QA pairs authored with Everllence engineers across four types (factual, table lookup, procedure synthesis, \emph{unanswerable}), with supporting-passage annotations.
    \item \textbf{Pipeline:} interchangeable chunkers and retrievers (BM25, multilingual dense, hybrid; optional reranker); generators: one Azure-hosted and one open-weight model, identical prompts; no fine-tuning.
    \item \textbf{Abstention:} prompt-level refusal with evidence-sufficiency self-check; measured as rejection rate (unanswerable) vs.\ false-refusal rate (answerable).
    \item \textbf{Evaluation:} retrieval recall@k/nDCG; answer correctness (expert rubric), groundedness (RAGAS \cite{es2023ragas}), citation accuracy; metric validation via agreement statistics between automatic scores and blinded expert ratings ($\sim$200 answers, 2--3 raters).
    \item \textbf{Scope/risk:} no agentic retrieval, no diagram understanding, no deployment. All RQs answerable on the public corpus alone if company access is delayed; LLM-judge fallback if expert capacity shrinks.
\end{itemize}

\section{Contribution and Value for Everllence}

(1) A quantitatively evaluated multilingual RAG assistant for industrial documentation; (2) an expert-validated measurement of abstention quality (negative rejection vs.\ false refusal) on real technical documents; (3) a released bilingual evaluation set and pipeline enabling regression testing. For Everllence: a deployment decision basis --- measured accuracy, groundedness, and refusal behaviour on its own document types and Azure infrastructure.

\section{Timeline}

{\small
\begin{tabularx}{\textwidth}{|p{1.6cm}|X|}
\hline
\rowcolor{lightgray} \textbf{Month} & \textbf{Work} \\
\hline
1 & Literature review; twin-corpus collection; table-aware ingestion; question typology. \\
\hline
2 & Evaluation-set authoring (workshop 1); baseline pipeline end-to-end. \\
\hline
3 & Retrieval experiments (RQ1); open-weight generator comparison. \\
\hline
4 & Abstention implementation and evaluation (RQ2); error analysis by type and language. \\
\hline
5 & RAGAS runs; blinded expert rating (workshop 2); metric validation (RQ3); company-corpus case study. \\
\hline
6 & Thesis writing, revision, release, submission. \\
\hline
\end{tabularx}}

\section{References}

{\footnotesize
\renewcommand{\refname}{}\vspace{-3.2em}
\begin{multicols}{2}
\raggedright
\begin{thebibliography}{9}\itemsep0pt

\bibitem{lewis2020rag}
P.\ Lewis et al.
\textit{Retrieval-Augmented Generation for Knowledge-Intensive NLP Tasks}. NeurIPS, 2020.

\bibitem{ji2023survey}
Z.\ Ji et al.
\textit{Survey of Hallucination in Natural Language Generation}. \textit{ACM Computing Surveys}, 55(12), 2023.

\bibitem{es2023ragas}
S.\ Es et al.
\textit{RAGAS: Automated Evaluation of Retrieval Augmented Generation}. EACL Demos, 2024.

\bibitem{saadfalcon2024ares}
J.\ Saad-Falcon, O.\ Khattab, C.\ Potts, M.\ Zaharia.
\textit{ARES: An Automated Evaluation Framework for RAG Systems}. NAACL, 2024.

\bibitem{chen2024benchmark}
J.\ Chen, H.\ Lin, X.\ Han, L.\ Sun.
\textit{Benchmarking Large Language Models in Retrieval-Augmented Generation}. AAAI, 2024.

\bibitem{karpukhin2020dpr}
V.\ Karpukhin et al.
\textit{Dense Passage Retrieval for Open-Domain Question Answering}. EMNLP, 2020.

\bibitem{robertson2009bm25}
S.\ Robertson, H.\ Zaragoza.
\textit{The Probabilistic Relevance Framework: BM25 and Beyond}. \textit{FnTIR}, 3(4), 2009.

\bibitem{everllence2025}
Everllence SE. \textit{Company and Product Information}. 2025. \url{https://www.everllence.com/}

\end{thebibliography}
\end{multicols}}

\vspace{0.4em}
\noindent
{\small
\begin{tabularx}{\textwidth}{X X}
\textbf{Student Signature:} & \textbf{Supervisor Signature:} \\[1.4em]
\rule{0.42\textwidth}{0.4pt} & \rule{0.42\textwidth}{0.4pt} \\
\end{tabularx}}

\end{document}
